Text detoxification can reduce abusive language, but it can also erase justified anger,
political critique, identity-related expression, or culturally specific rhetoric. We
therefore distinguish harmful insult reduction from disagreement or stance removal and
require human review for claims about non-toxicity and acceptability. The corpus
contains sensitive and offensive text. We limit quotation of examples, provide clear
annotation instructions and opt-out procedures, and avoid presenting automatic proxy
scores as safety outcomes. These methods should not be used to censor speakers or to
treat lower-resource language varieties as inherently more toxic or less fluent.
